\documentclass[conference]{IEEEtran}
\usepackage[utf8]{inputenc}
\usepackage[T1]{fontenc}
\usepackage{graphicx,booktabs,cite}
\usepackage[dvipsnames]{xcolor}
\usepackage{csquotes}
\usepackage{xspace}
\usepackage{acronym}
\usepackage{enumitem}
\usepackage{url}
\usepackage[hidelinks]{hyperref}
\graphicspath{{figures/}{../user-guide/screenshots/}}

\renewcommand{\sectionautorefname}{Section}
\renewcommand{\subsectionautorefname}{Section}
\renewcommand{\subsubsectionautorefname}{Section}
\renewcommand{\figureautorefname}{Figure}
\renewcommand{\tableautorefname}{Table}

% Draft notes. Remove all uses before submission.
\newcommand{\todo}[1]{\textcolor{red}{[TODO: #1]}}

% !TeX root = main.tex
% \acrodef defines the acronyms without printing a list.
\acrodef{API}{application programming interface}
\acrodef{CI}{continuous integration}
\acrodef{CPU}{central processing unit}
\acrodef{GPU}{graphics processing unit}
\acrodef{IPC}{inter-process communication}
\acrodef{LLM}{large language model}
\acrodef{PRISMA}{Preferred Reporting Items for Systematic Reviews and Meta-Analyses}
\acrodef{SE}{software engineering}
\acrodef{SLR}{systematic literature review}
\acrodef{SSRF}{server-side request forgery}

\newcommand{\API}{\ac{API}\xspace}
\newcommand{\APIs}{\acp{API}\xspace}
\newcommand{\CI}{\ac{CI}\xspace}
\newcommand{\CPU}{\ac{CPU}\xspace}
\newcommand{\GPU}{\ac{GPU}\xspace}
\newcommand{\GPUs}{\acp{GPU}\xspace}
\newcommand{\IPC}{\ac{IPC}\xspace}
\newcommand{\LLM}{\ac{LLM}\xspace}
\newcommand{\LLMs}{\acp{LLM}\xspace}
\newcommand{\PRISMA}{\ac{PRISMA}\xspace}
\newcommand{\SE}{\ac{SE}\xspace}
\newcommand{\SLR}{\ac{SLR}\xspace}
\newcommand{\SLRs}{\acp{SLR}\xspace}
\newcommand{\SSRF}{\ac{SSRF}\xspace}

% Tool name
\newcommand{\SaiLoR}{\textsc{SaiLoR}\xspace}

% !TeX root = main.tex
% Screenshot files used in the report, in one place so that they are easy to swap.
% The graphics path in main.tex searches figures/ first and then ../user-guide/screenshots/.
% All current screenshots come from the user guide (commit 1c0ec0f, 2026-09-07) and show the app before the move to Tauri.
% TODO: Retake every screenshot on the Tauri build of 2.0.0, store it in figures/, and change the names below only.
\newcommand{\shotAnnotate}{annotate-overview}
\newcommand{\shotSchema}{project-editor-schema}
\newcommand{\shotScreening}{screening-record}
\newcommand{\shotCompare}{consolidation-compare-popup}
\newcommand{\shotMerge}{git-merge-conflicts}


\title{\SaiLoR: Systematic Tool Support\\for Systematic Literature Reviews}
% TODO: final author list, affiliations and e-mail addresses.
\author{\IEEEauthorblockN{Jan Keim \todo{co-authors}}
\IEEEauthorblockA{Karlsruhe Institute of Technology (KIT)\\Karlsruhe, Germany\\\todo{e-mail}}}

\begin{document}
\maketitle

\begin{abstract}
%% Context
Systematic literature reviews are a common research method in software engineering.
They follow a protocol so that other researchers can trace and repeat each decision.
%% Problem
However, the review data is hard to keep traceable.
Reviewers record screening decisions and extracted values in spreadsheets, shared documents, or hosted services.
There, the link between an extracted value and its evidence in the paper can get lost.
Also, it is hard to reconstruct who decided what and when.
%% Idea
To tackle this, we present \SaiLoR, an open-source desktop tool that supports screening, data extraction, and consolidation of a review in one project.
The project is a set of plain JSON files that a schema describes and that git versions.
Several reviewers work independently in their own seats, and a consolidator resolves their disagreements with agreement statistics.
When reviewers synchronize through git, the tool merges concurrent changes field by field.
Optional large language model support proposes answers with evidence, and a human reviews every proposal.
%% Method
We describe the tool from the reviewer's view and explain its main architecture decisions, including the move from Electron to Tauri.
We also report how we keep the tool itself traceable.
%% Results
We documented 388 requirements and linked them to code, tests, and documentation.
Of these requirements, 92\% link to implementing source code and 90\% link to at least one automated test.
The test suites hold 2828 unit tests, 28 integration tests, and 100 Rust tests.
We release \SaiLoR 2.0.0 under the GPL-3.0 license and archive it on Zenodo.
\end{abstract}

\begin{IEEEkeywords}
Systematic Literature Review, Tool Support, Data Extraction, Requirements Traceability, Large Language Models
\end{IEEEkeywords}

% !TeX root = ../main.tex
\section{Introduction}\label{sec:introduction}

% Why SLRs matter
\SLRs are an established research method in \SE.
Kitchenham and Charters~\cite{kitchenham2007guidelines} adapted the method from medicine, and later guidelines added systematic mapping studies and snowballing~\cite{petersen2015guidelines,wohlin2014guidelines}.
The core promise of an \SLR is that other researchers can follow each step.
Therefore, reviewers define a protocol up front, select papers with explicit criteria, and extract data with a fixed form.
Reporting guidelines such as \PRISMA~\cite{page2021prisma} then ask authors to disclose how many papers they found, screened, and excluded, and why.

% Problem
However, the data that an \SLR produces is hard to keep traceable.
Reviewers often record screening decisions and extracted values in spreadsheets, shared documents, or hosted services.
In such a setup, the link between an extracted value and the passage in the paper that supports it is often only in the reviewer's head.
Also, two or more reviewers usually extract the same papers independently, and someone has to compare and consolidate their answers.
With a shared spreadsheet, it is hard to see who changed which cell and why.
Hassler et al.~\cite{hassler2014outcomes} collected barriers to the \SLR process in a community workshop, and several of these barriers concern tool support.
Recently, \LLMs promise to take over parts of screening and extraction.
However, studies in \SE and other fields find that their output still needs human control~\cite{felizardo2024chatgpt,khraisha2024llm}.

% Gap
Several tools support parts of the process.
Mapping studies and surveys list many of them~\cite{marshall2013tools,marshall2014tools,harrison2020software}.
Many of these tools are hosted web applications that keep the review data on a server, for example, Covidence~\cite{covidence}, Rayyan~\cite{ouzzani2016rayyan}, and CADIMA~\cite{kohl2018cadima}.
Other tools focus on one activity, such as screening with active learning in ASReview~\cite{vandeschoot2021asreview}.
To the best of our knowledge, no tool combines local, plain-file review data with a version history, schema-driven extraction with evidence links, independent reviewers with a consolidation step, and \LLM proposals that a human must accept.

% Idea
In this paper, we present \SaiLoR, an open-source desktop tool that supports \SLRs in a systematic way.
The name contains the letters S, L, and R of \enquote{systematic literature review}.
The main idea is to treat review data like source code.
A review is a project of plain JSON files.
A schema describes what reviewers extract, and the tool validates every answer against it.
Each reviewer writes into their own files, and git records every change.
When reviewers pull each other's work, the tool merges their files field by field.
Optional \LLM assistance proposes answers with quotes from the paper, but the reviewer decides which proposal to apply, and the project records that an \LLM was involved.

% Running example
We use the example project that ships with \SaiLoR as a running example.
The project is a small review on \enquote{Code Search and Program Repair}.
A screening project holds eight candidate papers.
Five of them are included and carried into an extraction project with a nested schema.
\autoref{fig:annotate} shows this extraction project in the main view of \SaiLoR.

% Contributions
In summary, we make the following contributions:
\begin{itemize}
    \item We present \SaiLoR, a desktop tool that supports screening, data extraction, and consolidation of an \SLR with several reviewers.
    \item We describe the main architecture decisions behind \SaiLoR and their reasons, including local-first storage, git as collaboration layer, the move from Electron to Tauri, and the control of \LLM assistance.
    \item We report how we keep the tool traceable with 388 requirements that we linked to source code, tests, and documentation.
    \item We release \SaiLoR 2.0.0 as open source under the GPL-3.0 license on GitHub and Zenodo~\cite{sailor2026zenodo}.
\end{itemize}

% Outline
The remainder of the paper is structured as follows: In \autoref{sec:related-work}, we discuss related work.
In \autoref{sec:approach}, we explain which \SLR activities \SaiLoR supports and which design principles guide it.
\autoref{sec:tool} walks through the tool from the reviewer's view.
\autoref{sec:architecture} describes the main architecture decisions.
In \autoref{sec:requirements}, we present the requirements and their trace links.
We discuss limitations in \autoref{sec:discussion} and conclude in \autoref{sec:conclusion}.

\section{Background and Related Work}
% TODO outline
% - SLR process (Kitchenham guidelines, PRISMA).
% - Tool support: reference managers, SLR platforms, qualitative-coding tools, LLM assistants.
% - Gap: structured annotation, multi-reviewer consolidation, git-based collaboration, local-first.

\section{Approach}
% TODO outline
% - Phases supported: protocol, search import, screening, data extraction, consolidation, reporting.
% - Schema-driven annotation, evidence links to PDF passages.
% - Reviewer seats and agreement metrics.

% !TeX root = ../main.tex
\section{SaiLoR from the Reviewer's View}\label{sec:tool}

In this section, we walk through \SaiLoR along the activities of an \SLR.
We use the running example from \autoref{sec:approach-example}.
% TODO: All figures show the user interface before the move to Tauri (user guide, commit 1c0ec0f). Retake them on the 2.0.0 build, see figures.tex.

\begin{figure*}[t]
    \centering
    \includegraphics[width=0.85\textwidth]{\shotAnnotate}
    \caption{Main view of \SaiLoR with the paper list on the left, the PDF in the middle, and the extraction form of the running example for reviewer seat 1 on the right}
    \label{fig:annotate}
\end{figure*}

\subsection{Planning: Protocol and Schema}\label{sec:tool-planning}

% Project editor
Reviewers start in the project editor.
There, they set the project title, the number of reviewer seats, and whether the project is a screening project.
An optional review protocol records the research questions, search strings, databases, search date, and further notes.
The main part of the editor is the schema, as shown in \autoref{fig:schema}.
Each row defines one node with its name, type, minimum and maximum count, options, description, and visibility condition.
In the figure, the field \emph{Exclusion Reason} is only visible if the field \emph{Relevant} has the answer \enquote{No}.
Reviewers can reorder and nest rows by drag and drop.

\begin{figure}[t]
    \centering
    \includegraphics[width=\columnwidth]{\shotSchema}
    \caption{Schema editor with field types, counts, options, and visibility conditions}
    \label{fig:schema}
\end{figure}

\subsection{Conducting: Import and Screening}\label{sec:tool-screening}

% Import and duplicates
Reviewers import references from BibTeX, RIS, or CSL-JSON files, and they can import a folder of PDFs including its subfolders.
During import, \SaiLoR sorts each record into one of three classes: new, certain duplicate, or probable duplicate.
A certain duplicate has the same normalized DOI or title as a paper in the project.
A probable duplicate has a fuzzy title similarity of at least $0.90$, or the same base title and an author similarity of at least $0.50$.
The reviewer must decide on each probable duplicate.
If the reviewer merges two records, \SaiLoR only fills fields that are empty in the existing paper.

% Screening
In a screening project, each reviewer decides per paper whether to include it, to exclude it with a reason, or to leave it undecided.
\autoref{fig:screening} shows the screening view of the running example.
The view shows the abstract instead of the PDF, because reviewers usually screen titles and abstracts first.
If the abstract was extracted from the PDF and not imported, \SaiLoR shows a warning, because the extraction can be wrong.
Reviewers can screen with the keyboard: \texttt{I} includes, \texttt{E} excludes, \texttt{U} clears the decision, and the number keys exclude with the matching reason.
After a decision, \SaiLoR moves on to the next paper.
A progress line shows how many papers the reviewer screened, included, and excluded.
The screening summary lists these counts with one row per exclusion reason, which is the input for a \PRISMA flow diagram~\cite{page2021prisma}.
Validation reports excluded papers without a reason.

\begin{figure}[t]
    \centering
    \includegraphics[width=\columnwidth]{\shotScreening}
    \caption{Screening view with the abstract, the decision buttons, and the progress line}
    \label{fig:screening}
\end{figure}

% Carry over
When screening is done, reviewers create the extraction project from the screening project.
\SaiLoR carries over all included papers, never carries over excluded papers, and lets the reviewer choose whether to carry over undecided papers.
The new project records where it came from and how many papers were included, excluded, and carried over.

\subsection{Conducting: Data Extraction}\label{sec:tool-extraction}

% Extraction view
\autoref{fig:annotate} shows the main view during extraction.
The left pane lists the papers, the middle pane shows the PDF, and the right pane shows the form that \SaiLoR generates from the schema.
Repeatable groups such as \emph{Findings} get buttons to add and remove entries.
A dot next to each paper shows how complete the paper is for the current seat.
Reviewers mark a paper as finished explicitly.
Thus, \enquote{finished} is a decision of the reviewer and not just the absence of empty fields.

% Evidence
Reviewers can copy selected text from the PDF into a field.
They can also highlight passages in five colors and add notes.
Each reviewer seat has its own marks.
A reviewer can link a mark to one or more fields as evidence.
The link popover of a field lists the linked marks, and a click on a mark jumps to the passage in the PDF.
In this way, the project stores where in the paper an answer comes from.
Reviewers can export their marks as standard PDF annotations into a new PDF file.

\subsection{Conducting: Consolidation}\label{sec:tool-consolidation}

% Seats
In a project with several reviewers, every reviewer must pick a seat before editing.
A reviewer only sees the answers of their own seat.
The consolidator works in the additional consolidation seat.

% Alignment
Before \SaiLoR can compare answers, it must find out which entries of a repeatable group belong together.
For example, two reviewers can list the findings of a paper in different orders.
\SaiLoR matches the entries with a maximum-weight assignment over their pairwise similarity, and it keeps entries with a similarity of $0.5$ or lower apart.
The alignment is deterministic, so that all reviewers see the same result.
Before \SaiLoR compares two answers, it trims them, converts them to lowercase, and collapses whitespace.
If at least two reviewers answered a field and all answers are equal, \SaiLoR fills the consolidated value automatically and marks it as auto-filled.

% Compare and adopt
For the remaining fields, the consolidator opens a comparison, as shown in \autoref{fig:compare}.
The comparison lists the answers of all reviewers.
The consolidator adopts one answer or marks the answers as equivalent.
The disagreements of a paper or of the whole project can be exported as text, for example, to discuss them in a meeting.

\begin{figure}[t]
    \centering
    \includegraphics[width=\columnwidth]{\shotCompare}
    \caption{Comparison of two reviewers' answers in the consolidation seat}
    \label{fig:compare}
\end{figure}

% Agreement
\SaiLoR also computes the inter-rater agreement.
For exactly two reviewers, it computes Cohen's $\kappa$~\cite{cohen1960coefficient}, and for two or more reviewers Fleiss' $\kappa$~\cite{fleiss1971measuring}.
In both cases, it also computes Krippendorff's $\alpha$ for nominal data~\cite{krippendorff2018content}, which tolerates missing answers.
\SaiLoR shows the agreement overall and per field.
If a statistic is not defined, for example, because all reviewers gave the same answer to every item, \SaiLoR reports the case instead of a number.
In a screening project, the agreement only considers the decision.

\subsection{Collaboration with Git}\label{sec:tool-git}

% Commit and pull
\SaiLoR has a built-in git view that uses the git installation of the user.
Instead of a line diff of JSON files, the view lists each changed field with its previous and current value.
Reviewers select which field changes they commit.
A pull fetches the upstream branch and reports whether the project is up to date, can be fast-forwarded, or needs a merge.
\SaiLoR only pushes when the reviewer asks for it.
It also supports branches, a project-scoped history, and stashes of the project's own files.

% Merge
For a merge, \SaiLoR does not use the line-based merge of git.
It runs a three-way merge on the level of fields.
If only one side changed a field, the merge keeps that change.
The schema is merged node by node, and answers under a removed field are kept.
Only if both sides changed the same field to different values does the reviewer have to decide.
\autoref{fig:merge} shows the dialog for these conflicts.
If a merge leaves conflicts in files that do not belong to the project, \SaiLoR aborts it and restores the previous state.

\begin{figure}[t]
    \centering
    \includegraphics[width=\columnwidth]{\shotMerge}
    \caption{Field-level resolution of merge conflicts after a pull}
    \label{fig:merge}
\end{figure}

\subsection{\acs{LLM} Assistance}\label{sec:tool-llm}

% Providers and opt-in
\LLM assistance is off unless the project enables it.
Reviewers can connect nine provider \APIs, among them Anthropic, OpenAI, Google Gemini, and any OpenAI-compatible server.
Local models run, for example, in Ollama or a llama.cpp server~\cite{llamacpp}.
\SaiLoR stores \API keys in the keychain of the operating system.
Every \LLM window shows a notice that the feature is a preview and that each proposal must be checked.

% Extraction proposals
For extraction, the model reads the paper and proposes answers.
\SaiLoR validates every proposal against the schema.
A review table shows each proposal with the field, the value, a quote from the paper, and a confidence.
In agent mode, the model can search and read pages of the paper, query OpenAlex and Crossref, and fetch web pages.
\SaiLoR checks whether the quoted evidence occurs in the paper, and an optional judge model checks the proposals in a bounded loop.
Nothing changes in the project until the reviewer applies the selected rows.
Applying is one undo step and does not overwrite existing answers.
\SaiLoR marks each field that the model wrote until the reviewer interacts with the field.
For each apply, \SaiLoR records the provider, the model, and the time in the project, and the paper shows a disclosure line such as \enquote{AI-assisted (2$\times$)}.

% AI seat and screening
In a project with several reviewers, the last seat can be dedicated to the model.
The answers of this seat go through the same consolidation as the answers of the human reviewers, so that the agreement statistics also cover the model.
For screening, the model proposes a decision and a reason per paper.
\SaiLoR never pre-selects an exclusion, and it pre-selects an inclusion only if the confidence is at least $0.8$.

% !TeX root = ../main.tex
\section{Architecture Decisions}\label{sec:architecture}

In this section, we describe the main architecture decisions behind \SaiLoR.
For each decision, we state the context, the decision, and its consequences.

\subsection{Overview}\label{sec:architecture-overview}

% Overview
\SaiLoR is a desktop application built with Tauri~2~\cite{tauri2}.
The user interface is a React application in TypeScript that runs in the web view of the operating system.
It renders PDFs with pdf.js and keeps the project state in a single store.
A backend in Rust provides the native capabilities, such as file access, git, the keychain, and network requests.
The user interface reaches the backend only through commands that the backend registers.
At the time of writing, the TypeScript code base has about 50,000 lines of code without tests and about 44,000 lines of tests.
The Rust backend has about 11,500 lines in 17 modules.

\subsection{Local Files Instead of a Server}\label{sec:architecture-local}

% Local-first
Hosted tools make collaboration easy, but they keep the review data on a server (cf.~\autoref{sec:related-work-tools}).
For a research tool, this setup has drawbacks.
The data can be hard to export completely, the service can change or end, and reviewers might not be allowed to upload papers.
Therefore, we decided to keep all data in local files, following the idea of local-first software~\cite{kleppmann2019local}.
The project file and the annotation files are JSON files that people can read.
\SaiLoR keeps unknown keys in the project file, so that an older version does not delete data of a newer version.
Each annotation file names the schema version it was written under, and validation reports files with an unknown schema version.
The serialization is deterministic, and opening a project writes nothing.
As a consequence, a version control system only sees changes that a reviewer made.
The price of this decision is that \SaiLoR needs its own collaboration mechanism.

\subsection{Git as Collaboration Layer}\label{sec:architecture-git}

% Git
We decided to use git for collaboration instead of building a synchronization server.
Git is common in \SE research, offers hosting at many institutions, and keeps a full history~\cite{ram2013git}.
To keep merges small, \SaiLoR splits the annotations into one file per paper and seat.
Two reviewers who work on different seats never edit the same annotation file.

% Field-level merge
However, a line-based merge does not understand JSON.
When two people edit the same file, for example, the consolidator and a reviewer who changes the schema, git can report conflicts in lines that do not conflict in their meaning.
Therefore, \SaiLoR runs its own three-way merge on the parsed files.
It compares both sides to the common base field by field.
A side that did not change a value does not influence the result.
Repeatable entries, project settings, and schema nodes are merged separately, and answers under a removed schema field are kept.

% Safety
Using the git installation of the user brings risks.
A cloned repository can contain a configuration that runs code, for example, hooks or a file system monitor.
\SaiLoR runs git only in the backend and overrides these settings on every call, including the pager, the editor, and the \texttt{ext} protocol.
It disables interactive prompts, limits the run time of each command, and validates clone URLs, paths, and branch names.
Git operations are only allowed on repositories that the user opened or cloned in the current session.
Before \SaiLoR commits or discards selected field changes, it checks that the file on disk still matches the reviewed state, and it refuses the write otherwise.

\subsection{From Electron to Tauri}\label{sec:architecture-tauri}

% Context
Up to version 1.x, \SaiLoR was an Electron application~\cite{electron}.
The Electron main process had about 5,600 lines of TypeScript and 98 \IPC handlers.
Electron bundles its own Chromium and Node.js runtime.
As a result, the macOS installer for arm64 had a size of 93~MB, and the main process could use the full Node.js \API.

% Decision
For version 2.0.0, we moved \SaiLoR to Tauri~2~\cite{tauri2}.
Tauri uses the web view of the operating system and a backend in Rust.
In a spike, a minimal Tauri shell produced installers of about 2~MB for macOS, Windows, and the Debian and RPM packages.
The AppImage for Linux has about 81~MB, because it bundles WebKitGTK.
The shell alone is thus about 2\% of the size of the old Electron installer.
The installers of the real application are larger because they include the user interface and its assets.
% TODO: Add the final installer sizes of the 2.0.0 release per platform.
\todo{final installer sizes of 2.0.0 per platform}

% Split
Size was not the only reason.
Tauri moves the security boundary into Rust and offers a capability system and a signed updater.
We split the responsibilities as follows.
Rust owns everything that needs authorization: file input and output, process spawning, downloads, and access to the keychain.
TypeScript keeps the parsing of files, the merge logic, and the assembly of exported PDFs, because this code was already tested and does not need native access.
\SaiLoR does not expose the generic file system and shell plugins of Tauri to the user interface.
The capability file grants the web view only the default core permissions.

% Consequences
The move has consequences for users and for us.
\SaiLoR now depends on three different web views: WKWebView on macOS, WebView2 on Windows, and WebKitGTK on Linux.
We checked that pdf.js renders on all three.
The minimum macOS version is 15.
Users of version 1.x have to install version 2.0.0 manually and enter their \API keys again.
Also, not all features of the Electron version were ported yet (cf.~\autoref{sec:discussion}).

\subsection{Security Model}\label{sec:architecture-security}

% Security
A review tool opens files and repositories from other people, and it can send data to \LLM providers.
Therefore, we treat the user interface as untrusted and enforce all checks in the backend.
\SaiLoR only saves a project to a path that the user opened or chose in a native dialog during the current session.
It serves a PDF only if the file lies inside the project folder.
For a PDF outside the folder, \SaiLoR asks the user and allows only this exact path for the current session.
Exports only write to destinations that the user picked in a save dialog, and \SaiLoR refuses symbolic links as targets.
Updates are signed, and \SaiLoR downloads and installs them only when the user asks for it.

% LLM security
For \LLM providers, \SaiLoR stores \API keys in the keychain and sends a key only to the origin of its provider.
It refuses redirects for requests that carry a key.
The agent mode can fetch web pages, which opens the door to \SSRF~\cite{owasp_ssrf}.
\SaiLoR only allows \texttt{http} and \texttt{https} URLs, resolves the host name, and refuses requests to loopback, private, and other non-public addresses.
It checks every redirect again, follows at most five redirects, stops after 20 seconds, reads at most 2~MB, and only accepts text content.
The fetched text is marked as untrusted data, and the system prompt tells the model to never follow instructions from it.

\subsection{Controlled \acs{LLM} Assistance}\label{sec:architecture-llm}

% Human in the loop
Prior studies found that \LLMs can help with screening and extraction but still make errors~\cite{felizardo2024chatgpt,syriani2024screening,khraisha2024llm}.
Therefore, we designed the \LLM support so that it cannot change the review data without a human decision.
Suggestions live outside the project until the reviewer applies them.
\SaiLoR validates them against the schema and checks quoted evidence against the text of the paper.
The durable record of each apply makes the use of \LLMs visible in the data, so that authors can report it.
Treating the model as an additional reviewer seat reuses the consolidation and agreement statistics, instead of adding a separate evaluation mechanism.

% Local runtime
Some reviewers cannot send papers to a hosted provider.
For them, \SaiLoR supports local models.
For chat models, users connect a local server such as Ollama, LM Studio, llama.cpp, or vLLM.
For small decision models, \SaiLoR can manage a llama.cpp server itself~\cite{llamacpp}.
It downloads the server and the models only from an allowlist of sources, verifies their SHA-256 checksums, and starts the server on \texttt{127.0.0.1} with a random key per session.
The server uses a \GPU when one is available and the \CPU otherwise.
We verified this runtime on macOS on arm64.
The managed installation needs a llama.cpp release with a specific route for decision models.
As long as no release includes this route, the setup dialog keeps the installation disabled and points users to a server that they run themselves.

\section{Requirements and Traceability}
% TODO outline
% - Requirement files and ID scheme (docs/requirements/index.md).
% - Links from requirements to code, tests, history and docs (docs/requirements/traceability.md).
% - Recovery method and coverage statistics.
% - How the links stay current (rule: update requirements with code changes).

\section{Discussion and Limitations}
% TODO outline
% - Threats to validity, scope limits, supported platforms.
% - Future work.

% !TeX root = ../main.tex
\section{Conclusion}\label{sec:conclusion}
\acresetall

In this paper, we looked into tool support that keeps the data of a \SLR traceable.
We presented \SaiLoR, an open-source desktop tool for screening, data extraction, and consolidation.
A \SaiLoR project consisted of plain JSON files that a schema described and that git versioned.
Reviewers worked independently in their own seats and linked their answers to evidence in the papers.
A consolidator resolved disagreements with the help of agreement statistics, and \SaiLoR merged concurrent changes field by field.
Optional \LLM assistance proposed answers with evidence, but a human reviewed every proposal before \SaiLoR applied it.

We also explained the main architecture decisions.
These were the local files instead of a server, git as collaboration layer, the security model, and the move from Electron to Tauri in version 2.0.0.
To keep the tool itself traceable, we documented 388 requirements.
Of these requirements, 92\% linked to implementing source code and 90\% linked to at least one automated test.
The test suites held 2828 unit tests, 28 integration tests, and 100 Rust tests.
We released \SaiLoR 2.0.0 under the GPL-3.0 license and archived it on Zenodo.

Regarding the tool, we want to port the remaining features of the Electron shell to Tauri and retest the user interface in the web views of all three operating systems.
We also want to verify the local model runtime on \GPUs with CUDA and Vulkan and enable the managed installation once llama.cpp provides the required route.

In future research, we want to evaluate \SaiLoR in a user study with reviewers who conduct a real \SLR.
We want to find out whether the tool reduces the effort of extraction and consolidation and whether the recorded history helps others to trace the review.
We also want to measure how accurate the \LLM proposals are and how they influence the decisions of reviewers.

Lastly, we want to have the requirements reviewed by people other than the developer and check the trace links automatically in the \CI.


\bibliographystyle{IEEEtran}
\bibliography{references}
\end{document}
